\documentclass{article}
\usepackage{lmodern}
\usepackage{amsmath}
\usepackage{listings}
\usepackage{fullpage}
\usepackage{parskip}
\usepackage{xcolor}
\lstset{
	basicstyle=\ttfamily,
	columns=fullflexible,
	frame=single,
	breaklines=true,
	postbreak=\mbox{\textcolor{red}{$\hookrightarrow$}\space},
}
\begin{document}
\title{Experiment `p\_2\_4\_a\_sum' Results}
\author{\today}
\date{}
\maketitle
\textbf{Experiment outcome: }SUCCESS
\\\textbf{Bad responses: }0
\\\textbf{Responses containing}~\texttt{assume}~\textbf{: }0
\\\textbf{Responses containing}~\texttt{expect}~\textbf{: }0
\\\textbf{Resolution attempts: }1
\\\textbf{Hard fails (resolution): }0
\\\textbf{Soft fails (resolution): }0
\\\textbf{Verification attempts: }1
\section*{Problem Specification}
\textbf{Problem name: }p\_2\_4\_a\_sum
\\\textbf{Natural language statement: }Write a method that takes two integers and then returns the sum.
\\\textbf{Method signature: }p\_2\_4\_a\_sum(x: int, y: int) returns (result: int)
\subsection*{Ensures}
\begin{itemize}
\item \texttt{result == x + y}
\end{itemize}
\clearpage
\section*{GenAI interactions}
Below you will find all interactions between the `user' (program) and the `assistant' (OpenAI).
\subsection*{Program $\rightarrow$ GenAI}
\begin{lstlisting}
You are given the following task to perform in Dafny:

Write a method that takes two integers and then returns the sum.

The signature should be:

p_2_4_a_sum(x: int, y: int) returns (result: int)

The method should respect the following contract:

ensures result == x + y

Produce and show only the Dafny body of this method, including the curly braces that surround it. Do not show the signature nor contract. You must not use 'assume' nor 'expect' anywhere in your code.
\end{lstlisting}
\subsection*{GenAI $\rightarrow$ Program}
\textbf{System fingerprint: }null
\\\textbf{ID: }chatcmpl-DmdWxY2idkpPNsFIIl9RLeAllR2uU
\\\textbf{Model: }gpt-5.5-2026-04-23
\\\textbf{Created at: }1780484071
\\\textbf{Finish reason: }stop
\\\textbf{Usage: }260 tokens in, and 56 tokens out
\begin{lstlisting}
{
  result := x + y;
}
\end{lstlisting}
\section*{Final Program}
\begin{lstlisting}
method p_2_4_a_sum(x: int, y: int) returns (result: int)
	ensures result == x + y
{
  result := x + y;
}
\end{lstlisting}
\section*{Total Token Usage}
\textbf{Input tokens: }260
\\\textbf{Output tokens: }56
\\\textbf{Reasoning tokens: }28
\\\textbf{Sum of `total tokens': }316
\section*{Experiment Timings}
\textbf{Overall Experiment} started at 1780484071592, ended at 1780484077109, lasting 5517ms (5.52 seconds)
\\\textbf{Iteration \#1} started at 1780484071593, ended at 1780484077108, lasting 5515ms (5.52 seconds)
\end{document}
